% Acronyms
% \newacronym{api}{API}{Application Program Interface}
% \newacronym{uml}{UML}{Unified Modeling Language}

% Glossary entries
\newglossaryentry{apig}{
    name={API},
    description={in informatica con il termine \emph{Application Programming Interface API} (ing. interfaccia di programmazione di un'applicazione) si indica ogni insieme di procedure disponibili al programmatore, di solito raggruppate a formare un set di strumenti specifici per l'espletamento di un determinato compito all'interno di un certo programma. La finalità è ottenere un'astrazione, di solito tra l'hardware e il programmatore o tra software a basso e quello ad alto livello semplificando così il lavoro di programmazione}
}

\newglossaryentry{umlg}{
    name={UML},
    description={in ingegneria del software \emph{UML, Unified Modeling Language} (ing. linguaggio di modellazione unificato) è un linguaggio di modellazione e specifica basato sul paradigma object-oriented. L'\emph{UML} svolge un'importantissima funzione di ``lingua franca'' nella comunità della progettazione e programmazione a oggetti. Gran parte della letteratura di settore usa tale linguaggio per descrivere soluzioni analitiche e progettuali in modo sintetico e comprensibile a un vasto pubblico}
}

\newglossaryentry{mcpg}{
    name={MCP},
    description={non ancora implementato}
}

\newglossaryentry{toolg}{
    name={tool},
    text={Tool},
    description={non ancora implementato}
}

\newglossaryentry{htmlg}{
    name={HTML},
    description={non ancora implementato}
}

\newglossaryentry{jsong}{
    name={JSON},
    description={non ancora implementato}
}

\newglossaryentry{dockerg}{
    name={Docker},
    description={non ancora implementato}
}

\newglossaryentry{crawlingg}{
    name={Web crawling},
    description={non ancora implementato}
}

\newglossaryentry{jwtg}{
    name={JWT},
    description={non ancora implementato}
}

\newglossaryentry{prompt-engineeringg}{
    name={Prompt engineering},
    description={non ancora implementato}
}

\newglossaryentry{llmg}{
    name={LLM},
    description={non ancora implementato}
}

\newglossaryentry{cachingg}{
    name={caching},
    description={non ancora implementato}
}

\newglossaryentry{frameworkg}{
    name={framework},
    description={non ancora implementato}
}

\newglossaryentry{sseg}{
    name={SSE},
    description={non ancora implementato}
}

\newglossaryentry{httpg}{
    name={HTTP},
    description={non ancora implementato}
}

\newglossaryentry{websocketg}{
    name={WebSocket},
    description={non ancora implementato}
}

\newglossaryentry{mvvmg}{
    name={MVVM},
    description={non ancora implementato}
}